\subsection{Transformations (Klein-Lie, 1869-1872).}


Around 1860, the theory of permutation groups of a finite set is developing and begins to be used (Serret, Kronecker, Mathieu, Jordan). On the other hand, the theory of invariants, then in full flight, familiarises mathematicians with certain infinite sets of geometric transformations closed under composition (notably linear or projective transformations). But, before the work of 1868 of Jordan {[}174 b{]} on the ``movement groups'' (closed subgroups of the groups of displacements in Euclidean space of 3 dimensions), it does not seem as if any conscious link had been established between these two streams of ideas.

In 1869, the young Felix Klein (1849-1925), a pupil of Plücker, becomes the friend at Berlin of the Norwegian Sophus Lie (1842-1899), a few years older, being brought together by their common interest in the ``geometry of straight lines'' of Plücker and notably the theory of complexes of straight lines (p.~133). It is around this period that Lie conceives of one of his most original ideas, the introduction of the notion of invariant in Analysis and in differential geometry; one of the sources of this is his observation that the classical methods of integration ``by quadratures'' of differential equations all rely on the fact that the equation is invariant for a ``continuous'' family of transformations. It is from 1869 that is dated the first work (drawn up by Klein) where Lie uses this idea; he studies there the ``Reye complex'' (a set of straight lines intersecting the faces of a tetrahedron in 4 points having a given cross-ratio) and the curves and surfaces having as tangents straight lines of this complex ({[}202{]}, vol.~I, Abh. V, pp.~68-77): his method relies on the invariance of the Reye complex under the commutative group with 3 parameters (a maximal torus of PGL(4, C)) leaving invariant the corners of the tetrahedron. This same idea dominates the work written together by Klein and Lie when they found themselves in Paris in the spring of 1870 ({[}182{]}, v. I, pp.~416-420); they essentially determine there the connected commutative subgroups of the projective group of the plane PGL(3, C), and study the geometric properties of their orbits (under the name of V curves or surfaces); that gives them, by a uniform procedure, properties of varied curves, algebraic or transcendental, such as \(y = cx^m\) or the logarithmic spirals. Their testimony agrees in underlining the profound impression that was produced on them by the theories of Galois and of Jordan (the commentary of Jordan on Galois had appeared in Math. Annalen in 1869; in any case, Lie had heard of the theory of Galois already in 1863). Klein, who in 1871, begins to be interested in non-Euclidean geometries, sees there the beginning of his research for a classification principle for all known geometries, research that was to lead him in 1872 to the ``Erlangen programme''. On his side, Lie in a letter of 1873 to A. Mayer ({[}202{]}, vol.~5, p.~584), dates from his stay in Paris the origin of his ideas on transformations groups, and in a work of 1871 ({[}202{]}, vol.~I, Abh. XII, pp.~153-214), he uses already the term ``transformation group'' and states explicitly the problem of the determination of all the subgroups (``continuous or discontinuous'') of GL(n, C). Truth to tell, Klein and Lie must both have had difficulty in penetrating this new mathematical universe, and Klein speaks of the ``Treatise'' of Jordan, newly appeared, as a ``book sealed with seven seals'' ({[}182{]}, v. I, p.~51); he writes elsewhere in connection with ({[}182{]}, v. I, pp.~424-459): ``It is to Lie that belongs all that is concerned with the heuristic idea of a continuous group of operators, in particular all that touches on the integration of differential equations or partial differential equations. All the notions that he developed later in his theory of continuous groups were already to be found in embryo with him, but so little elaborated however, that I had to convince him of many details, for example at the beginning even the existence of the curves V, during the course of long conversations'' ({[}182{]}, v. I, p.~415).

